% Appendix I -- the doubling graph, certified: a machine-checked theorem and the ideas behind it.
\section{\appendixItitle}\label{app:lean}

This appendix carries the mathematics of Appendix~\ref{app:doubling} --- the doubling graph
\eqref{eq:counterexample_graph} and the family of backward policies it carries --- with the
proofs delegated to a machine-checked certificate rather than written out in natural language,
and it is meant to replace that appendix rather than to stand beside it.
Theorem~\ref{theo:lean_main} is the certified statement; it is narrower than
Theorem~\ref{theo:doubling_main} in six places, which \S\ref{sec:lean_scope} names, with
where each is proved instead when it is proved at all. \S\ref{sec:lean_model} quotes the formal model the
theorem quantifies over, \S\ref{sec:lean_certificate} the fifteen Lean theorems that carry it
and the axioms they rest on, and \S\ref{sec:lean_ideas} says, in prose and without proving
anything, how each is proved. The numbers measured on truncations are not repeated here: they
stay in \S\ref{sec:doubling_measured} and \S\ref{sec:doubling_cost}, which this appendix does
not replace.

\begin{definition}[standing notation and hypotheses of this appendix]\label{def:lean_setting}
The graph is \eqref{eq:counterexample_graph}, carrying the source $s_0$, the ladder
$1,2,3,\dots$, the sink $s_f$ and the forward edges $s_0\rightarrow1$, $j\rightarrow j+1$,
$2j\rightarrow j$ and $j\rightarrow s_f$; the predecessors of a ladder state $j$ are $j-1$ and
$2j$, so the backward chain at $j$ either \emph{decrements} to $j-1$ or \emph{doubles} to $2j$,
and the \emph{loop closure} adds the wrap $s_f\rightarrow s_0$. The backward policy is
\begin{equation}\label{eq:lean_policy}
    \BackwardPolicy(j\rightarrow2j)=\varepsilon(j),\qquad
    \BackwardPolicy(j\rightarrow j-1)=1-\varepsilon(j)\quad(j\geq1),\qquad
    \BackwardPolicy(s_0\rightarrow s_f)=1,
\end{equation}
with $j-1=0$ read as $s_0$, and the \emph{target row}
$\widehat\BackwardPolicy(s_f\rightarrow\cdot)$ is a probability supported in $\{1,\dots,d\}$ of
mean $\bar\jmath$, the integer $d\geq1$ being a parameter of the setting and $\bar\jmath$
the mean of its row. The family studied throughout is
\begin{equation}\label{eq:lean_family}
    \varepsilon_{c,s}(j):=\frac{c}{(j+1)^s},
\end{equation}
and its \emph{standing range} is $s\geq0$, $0<c<2^s$, on which $\varepsilon_{c,s}$ takes its
values in $(0,1)$ with supremum $\varepsilon_{\max}=c\,2^{-s}<1$.

Write $(X_n)_{n\geq0}$ for the backward chain of the loop closure, $\sigma$ for its hitting
time of $s_0$, and $\bar\sigma$ for the expected backward-trajectory length of
\citet{morozov2025revisiting}, the average of $\mathbb E(\sigma\mid X_0=j)$ over the target
row. When the loop closure carries an invariant probability, $\lambda$ denotes one,
$\lambda_j:=\lambda(\{j\})$ the mass of a ladder state, $L(m):=\lambda([m,+\infty))$ the tail,
and $f_m:=\mathbf 1_{[m,+\infty)}-L(m)$, for $m>d$, the \emph{centred tail indicator} at the
cut $m$. The operator $P_\star$ is the \emph{function action} of the backward kernel,
$(P_\star f)(x)=\mathbb E(f(X_1)\mid X_0=x)$, defined at every function on the states and a
contraction of $L^p(\lambda)$ for every $p\in[1,+\infty]$; $P$ is the density action of that
kernel on $L^2(\lambda)$; $\Pi$ is the $\lambda$-mean projection
$f\mapsto\bigl(\int f\,d\lambda\bigr)\mathbf 1$ and $\ker\Pi$ the functions of zero
$\lambda$-mean in $L^2(\lambda)$; $\widehat\beta_n:=\|P^n-\Pi\|_{L^2(\lambda)}$; and
\begin{equation}\label{eq:lean_Bhat}
    S:=\left(\mathrm{Id}-P_\star+\Pi\right)^{-1}-\Pi,\qquad
    \widehat B:=\|S\|_{L^2(\lambda)},
\end{equation}
with the convention $\widehat B:=+\infty$ when no bounded operator $R$ on $L^2(\lambda)$
satisfies $R(\mathrm{Id}-P_\star)=\mathrm{Id}-\Pi$. The \emph{Cram\'er equation} of the family
is
\begin{equation}\label{eq:lean_cramer}
    c\left(2^p-1\right)\;=\;p ,
\end{equation}
and $p_*=p_*(c)$ denotes its non-zero real root. For an even integer $K\geq d$, the
\emph{truncation at $K$} is the chain obtained from the loop closure by deleting every doubling
edge $j\rightarrow2j$ with $2j>K$, the decrement out of such a state carrying probability $1$;
$\lambda^K$ denotes an invariant probability of it and $\widehat B_K$ its constant
\eqref{eq:lean_Bhat}. Definition~\ref{def:doubling_setting} carries the same objects
at length.
\end{definition}

\begin{theorem}[the doubling graph, in the form the certificate carries]\label{theo:lean_main}
In the setting of Definition~\ref{def:lean_setting}, with $(c,s)$ in the standing range:
\begin{enumerate}
    \item \emph{(irreducibility and the phase diagram)} the loop closure is irreducible and
    every invariant probability of it is positive at every state; and, for $c>0$, it carries no
    invariant probability at $s<1$, at least one at $s=1$ and $0<c<1$, none at $s=1$ and
    $1\leq c$, and at least one at $s>1$;
    \item \emph{(the invariant measure at $s=1$)} at $s=1$ and $0<c<1$, with $\lambda$ an
    invariant probability and $p_*>1$ the root of \eqref{eq:lean_cramer}, there is $C>0$ with
    $\lambda_m m^{p_*}\rightarrow C$ and $(p_*-1)m^{p_*-1}L(m)\rightarrow C$ --- that is
    $L(m)=\frac C{p_*-1}m^{1-p_*}\bigl(1+o(1)\bigr)$ --- and there are $\vartheta\in(0,1)$ and
    $c_6>0$ with $\left|\lambda_m m^{p_*}-C\right|\leq c_6m^{-\vartheta}$ for every $m\geq1$;
    the constant is determined by the boundary data in two senses: two invariant probabilities
    of the loop closure that agree at the ladder states $1,\dots,d$ agree along the whole
    ladder; and there are $\nu_1,\dots,\nu_d\geq0$ with $\sum_j\nu_j>0$ and $\nu_j=0$
    whenever $2j\leq d$ such that every invariant probability $\lambda'$ of the loop closure
    has $\lambda'_m m^{p_*}\rightarrow\sum_{j\leq d}\nu_j\lambda'_j$;
    \item \emph{(the diffusion operator)} at every $(c,s)$ of the standing range and every
    invariant probability $\lambda$: for every $p\in[1,+\infty)$ the centred tail indicators
    drive the $L^p(\lambda)$ Rayleigh quotient to zero --- for every $t>0$ some $m>d$ has
    $\|(\mathrm{Id}-P_\star)f_m\|^p_{L^p(\lambda)}\leq t\|f_m\|^p_{L^p(\lambda)}$ --- and no
    constant $B$ satisfies $\|f\|^p_{L^p(\lambda)}\leq B\|(\mathrm{Id}-P_\star)f\|^p_{L^p(\lambda)}$
    for every bounded $f$ of zero $\lambda$-mean; at $p=+\infty$, no constant $B$ satisfies
    $M\leq B\,D$ whenever a bounded $f$ of zero mean attains $|f|\geq M$ somewhere and has
    $|(\mathrm{Id}-P_\star)f|\leq D$ everywhere; and at $p=2$, no bounded operator $R$ on
    $L^2(\lambda)$ satisfies $R(\mathrm{Id}-P_\star)=\mathrm{Id}-\Pi$, so $\widehat B=+\infty$,
    the partial sums of $\sum_{n\geq0}(P_\star^n-\Pi)$ do not converge in operator norm, and
    $\sum_{n\geq0}\widehat\beta_n=+\infty$. At $s=1$ and $0<c<1$ the expected
    backward-trajectory length is finite, $\bar\sigma=\bar\jmath/(1-c)$, and there are
    $0<c_1\leq c_2$ and an integer $m_2>d$ with
    \begin{equation}\label{eq:lean_rate}
        c_1\,m\,\left\|\left(\mathrm{Id}-P_\star\right)f_m\right\|^2_{L^2(\lambda)}
        \;\leq\;4\,c_2\left(p_*-1\right)\left\|f_m\right\|^2_{L^2(\lambda)}
    \end{equation}
    for every $m\geq m_2$ --- a $\sqrt m$ rate along the trial functions, and not a bound on
    $\widehat B$, which is $+\infty$ there;
    \item \emph{(an unsolvable flow-matching problem)} at every $(c,s)$ of the standing range
    and every invariant probability $\lambda$, the subspace
    $\left(\mathrm{Id}-P_\star\right)L^2(\lambda)$ is a dense proper subspace of $\ker\Pi$, and
    there are probability densities $\finit,\fterm\in L^2(\lambda)$ such that no
    $f\in L^2(\lambda)$ satisfies $\left(\mathrm{Id}-P_\star\right)f=\finit-\fterm$;
    \item \emph{(the truncation)} for every even $K\geq d$ the truncation at $K$ is irreducible
    on its states and every invariant probability of it is positive at every one of them; it
    carries exactly one such probability $\lambda^K$, and admits a bounded operator $R$ on
    $L^2(\lambda^K)$ with
    $R(\mathrm{Id}-P_\star)=(\mathrm{Id}-P_\star)R=\mathrm{Id}-\Pi$ and $\Pi R=R\Pi=0$, so
    that $\widehat B_K<+\infty$; and at $s=1$ and $0<c<1$ there are $c_8>0$ and an integer $K_0$
    such that $c_8K\leq\|R\|^2$ for every even $K\geq K_0$ and every bounded $R$ on
    $L^2(\lambda^K)$ with $R(\mathrm{Id}-P_\star)=\mathrm{Id}-\Pi$ --- that is
    $\widehat B_K\geq\sqrt{c_8}\sqrt K$.
\end{enumerate}
\end{theorem}

The two conditions come apart on this graph. That of \citet{morozov2025revisiting} is repaired
by every member with $s=1$ and $c<1$, which is what $\bar\sigma=\bar\jmath/(1-c)$ says; the
summable-mixing hypothesis of Theorem~\ref{theo:universality_L2} is repaired by no positive
recurrent member, which is item \emph{(3)} read at $p=2$. What item \emph{(4)} refutes is
\emph{exact} $L^2$ flow matching --- a pair of $L^2$ densities the loop closure cannot balance
at all --- and not the weak universality of Theorem~\ref{theo:universality_L2_full}, whose
conclusion is a defect infimum of zero: the range of $\mathrm{Id}-P_\star$ in $\ker\Pi$ being
dense is exactly why the weak form is untouched.

\subsection{The model}\label{sec:lean_model}

A machine-checked proof certifies that each proof proves its statement; what it cannot certify
is that the statement is the paper's. Four modelling decisions carry that burden here, and each
is legible off the source below.

Three names collide with the paper's, and a reader checking a quotation needs the dictionary
before the first of them. In the Lean, $\texttt{S}$ is a $\texttt{Setting}$ --- the data of the
graph and the policy --- where the paper's $S$ is the diffusion operator \eqref{eq:lean_Bhat},
which appears there under the name $\texttt{Sop}$; $\texttt{L}$ is an invariant probability and
$\texttt{L.lam}$ is the paper's $\lambda$, where the paper's $L(m)$ is the tail, which is
$\texttt{L.tailMass}$; and $\texttt{p}$ is the exponent the paper writes $p_*$.

The states are a ladder indexed by $\mathbb N$ and a sink, and \emph{the source $s_0$ is the
ladder index $0$}, which is Definition~\ref{def:lean_setting}'s own convention --- so the
decrement is uniformly $\texttt{lad}\ (j+1)\rightarrow\texttt{lad}\ j$ with no case at the foot
of the ladder. A single parameter $\texttt{cap}$ of type $\texttt{Option}\ \mathbb N$ carries
both chains: $\texttt{none}$ is the loop closure, where every doubling edge is present, and
$\texttt{some}\ K$ the truncation at $K$, where $\texttt{HasDouble}$ deletes the edge out of $j$
as soon as $2j>K$. The predicate $\texttt{OnChain}$ names the states the truncation lives on.
\begin{lstlisting}[style=lean]
inductive St where
  | lad : ℕ → St
  | sink : St
  deriving DecidableEq, Repr

def HasDouble : Option ℕ → ℕ → Prop
  | none, _ => True
  | some K, j => 2 * j ≤ K

def OnChain : Option ℕ → St → Prop
  | none, _ => True
  | some K, .lad j => j ≤ K
  | some _, .sink => True
\end{lstlisting}

The standing hypotheses of Definition~\ref{def:lean_setting}, minus the family, are one
structure. The target row is finitely supported by hypothesis, so every expression involving it
is a finite sum; $\varepsilon_{\max}$ is a field rather than a supremum, because only
$\varepsilon(j)\leq\varepsilon_{\max}<1$ is ever used; and $\bar\jmath$ is not a field at all
but is derived from the row, $\texttt{Setting.jbar}$. The family \eqref{eq:lean_family} and the
Cram\'er equation \eqref{eq:lean_cramer} are two functions. The structure is quoted with its
field docstrings elided, which is the one place a quotation below is not the whole block.
\begin{lstlisting}[style=lean]
structure Setting where
  d : ℕ
  d_pos : 1 ≤ d
  eps : ℕ → ℝ
  eps_pos : ∀ ⦃j⦄, 1 ≤ j → 0 < eps j
  epsMax : ℝ
  eps_le : ∀ ⦃j⦄, 1 ≤ j → eps j ≤ epsMax
  epsMax_lt_one : epsMax < 1
  row : ℕ → ℝ
  row_nonneg : ∀ j, 0 ≤ row j
  row_supp : ∀ ⦃j⦄, row j ≠ 0 → 1 ≤ j ∧ j ≤ d
  row_sum : ∑ j ∈ Finset.Icc 1 d, row j = 1

noncomputable def epsCS (c s : ℝ) (j : ℕ) : ℝ := c / ((j : ℝ) + 1) ^ s

noncomputable def psi (c t : ℝ) : ℝ := c * ((2 : ℝ) ^ t - 1) - t
\end{lstlisting}

The kernel enters chiefly through $P_\star$, and $P_\star$ is given \emph{in closed form as a
total function} on all of $\texttt{St}\rightarrow\mathbb R$ --- three cases for the three kinds
of state, and no integrability side condition anywhere. That is the second decision, and it is
what lets every statement at the finite exponents, and at $p=+\infty$, be an inequality between
real numbers rather than a statement about a measure. It does not carry everything: the density
action enters as well, in $\widehat\beta_n$, and the support of the kernel is encoded separately
by $\texttt{Edge}$, which is what the irreducibility statements read.
\begin{lstlisting}[style=lean]
noncomputable def pstar (S : Setting) (cap : Option ℕ) (f : St → ℝ) : St → ℝ
  | .lad 0 => f .sink
  | .lad (j + 1) =>
      if HasDouble cap (j + 1) then
        S.eps (j + 1) * f (.lad (2 * (j + 1))) + (1 - S.eps (j + 1)) * f (.lad j)
      else
        f (.lad j)
  | .sink => ∑ k ∈ Finset.Icc 1 S.d, S.row k * f (.lad k)
\end{lstlisting}

The third decision is that ``an invariant probability'' is a structure and not a theorem.
Invariance is carried in the integral form $\int P_\star f\,d\lambda=\int f\,d\lambda$ against
bounded test functions, which is the form every use consumes; positivity is carried on chain,
because one state type serves two chains and the truncation genuinely does not live on all of
it. Every statement of Theorem~\ref{theo:lean_main} that names $\lambda$ therefore quantifies
over such a structure rather than asserting that one exists --- existence is
$\texttt{main\_phase}$, and it is the hard half of item \emph{(1)}. Irreducibility is the
reflexive-transitive closure of the positive transitions of \eqref{eq:lean_policy}, and
$\texttt{PreStat}$ is $\texttt{Stat}$ without positivity, which is what makes ``irreducible
implies positive everywhere'' a theorem rather than an assumption.
\begin{lstlisting}[style=lean]
inductive Edge (S : Setting) (cap : Option ℕ) : St → St → Prop
  | wrap : Edge S cap (.lad 0) .sink
  | dec (j : ℕ) : Edge S cap (.lad (j + 1)) (.lad j)
  | dbl (j : ℕ) (h : HasDouble cap (j + 1)) : Edge S cap (.lad (j + 1)) (.lad (2 * (j + 1)))
  | row (k : ℕ) (hk : k ∈ Finset.Icc 1 S.d) (hpos : 0 < S.row k) : Edge S cap .sink (.lad k)

abbrev Reach (S : Setting) (cap : Option ℕ) : St → St → Prop :=
  Relation.ReflTransGen (Edge S cap)

structure Stat (S : Setting) (cap : Option ℕ) where
  lam : St → ℝ
  nonneg : ∀ x, 0 ≤ lam x
  pos : ∀ ⦃x⦄, OnChain cap x → 0 < lam x
  vanish : ∀ ⦃x⦄, ¬ OnChain cap x → lam x = 0
  summable : Summable lam
  total : ∑' x, lam x = 1
  inv : ∀ f : St → ℝ, (∃ C, ∀ x, |f x| ≤ C) →
          ∑' x, lam x * pstar S cap f x = ∑' x, lam x * f x

structure PreStat (S : Setting) (cap : Option ℕ) where
  lam : St → ℝ
  nonneg : ∀ x, 0 ≤ lam x
  vanish : ∀ ⦃x⦄, ¬ OnChain cap x → lam x = 0
  summable : Summable lam
  total : ∑' x, lam x = 1
  inv : ∀ f : St → ℝ, (∃ C, ∀ x, |f x| ≤ C) →
          ∑' x, lam x * pstar S cap f x = ∑' x, lam x * f x
\end{lstlisting}

The fourth decision is where the norms live. At finite $p$ the statements are in a real
\emph{mass} layer, $\texttt{L.mass}\ p\ f=\|f\|^p_{L^p(\lambda)}$, a plain sum of non-negative
reals; the tail $L(m)$ and the centred tail indicator $f_m$ are functions of the same kind, and
so are the growth condition on $\varepsilon$ along the powers of two, the recursion for
$\mathbb E(\sigma\wedge n\mid X_0=\cdot)$ and the truncated averages whose supremum is
$\bar\sigma$.
\begin{lstlisting}[style=lean]
noncomputable def Stat.mass (L : Stat S cap) (p : ℝ) (f : St → ℝ) : ℝ :=
  ∑' x, |f x| ^ p * L.lam x

noncomputable def Stat.tailMass (L : Stat S cap) (m : ℕ) : ℝ := ∑' x, L.lam x * tailInd cap m x

noncomputable def Stat.centredTail (L : Stat S cap) (m : ℕ) : St → ℝ :=
  fun x => tailInd cap m x - L.tailMass m

def GrowthCond (S : Setting) : Prop :=
  Tendsto (fun D : ℕ => (∑ i ∈ Finset.range D, Real.log (1 / S.eps (2 ^ i))) / 2 ^ D)
    atTop (𝓝 0)

noncomputable def hitExp (S : Setting) (cap : Option ℕ) : ℕ → St → ℝ
  | 0 => fun _ => 0
  | n + 1 => fun x =>
      match x with
      | .lad 0 => 0
      | .sink => 0
      | .lad (j + 1) => 1 + pstar S cap (hitExp S cap n) (.lad (j + 1))

noncomputable def sbar (S : Setting) (cap : Option ℕ) (n : ℕ) : ℝ :=
  ∑ k ∈ Finset.Icc 1 S.d, S.row k * hitExp S cap n (.lad k)
\end{lstlisting}

Only the $L^2$ operator statements --- items \emph{(3)} at $p=2$, \emph{(4)} and \emph{(5)} ---
need a genuine Hilbert space, and there they run on Mathlib's $L^p$ spaces over the measure
$\lambda$ read as a counting measure with density, with $P_\star$ and $\Pi$ continuous linear
maps on it, $\widehat\beta_n$ the norm of $P^n-\Pi$ for the density action, and $\ker\Pi$ the
orthogonal complement of the constants. That $P_\star$ is the $L^2(\lambda)$-adjoint of $P$ is a
theorem of the library, $\texttt{Stat.adjoint\_pstarL2}$, and not a convention; it is not among
the fifteen quoted below, and item \emph{(3)}'s clause on $\sum_n\widehat\beta_n$ is where the
theorem is consumed.
\begin{lstlisting}[style=lean]
noncomputable def mu : Measure St :=
  Measure.count.withDensity (fun x => ENNReal.ofReal (L.lam x))

noncomputable def pstarL2 (hrow : RowOnChain S cap) : Lp ℝ 2 L.mu →L[ℝ] Lp ℝ 2 L.mu :=
  LinearMap.mkContinuous
    { toFun := L.pstarLp
      map_add' := L.pstarLp_add hrow
      map_smul' := fun c F => L.pstarLp_smul hrow c F } 1
    fun F => by simpa using L.norm_pstarLp_le F

noncomputable def piL2 : Lp ℝ 2 L.mu →L[ℝ] Lp ℝ 2 L.mu :=
  (innerSL ℝ L.oneLp).smulRight L.oneLp

noncomputable def betaHat (hrow : RowOnChain S cap) (n : ℕ) : ℝ :=
  ‖(L.densL2 hrow) ^ n - L.piL2‖

noncomputable def partialSum (P Pi : E →L[ℝ] E) (N : ℕ) : E →L[ℝ] E := ∑ n ∈ Finset.range N, (P ^ n - Pi)

noncomputable def defect : Lp ℝ 2 L.mu →L[ℝ] Lp ℝ 2 L.mu := 1 - L.pstarL2 (rowOnChain_none S)

noncomputable def kerPi : Submodule ℝ (Lp ℝ 2 L.mu) := (ℝ ∙ L.oneLp)ᗮ

noncomputable def defectRange : Submodule ℝ (Lp ℝ 2 L.mu) :=
  LinearMap.range (L.defect : Lp ℝ 2 L.mu →ₗ[ℝ] Lp ℝ 2 L.mu)
\end{lstlisting}

\subsection{The certificate, and its axioms}\label{sec:lean_certificate}

Fifteen theorems carry Theorem~\ref{theo:lean_main}, all of them in the module
$\texttt{GFNBounds/Doubling/Main.lean}$ of the supplementary library, and each is quoted below
in full, statement only. Item \emph{(1)} is two: irreducibility with positivity, and the phase
diagram, each of whose four rows is the existence or the non-existence of an invariant
probability.
\begin{lstlisting}[style=lean]
theorem main_irreducible (P : PreStat S none) :
    (∀ x y : St, Reach S none x y) ∧ (∀ y : St, 0 < P.lam y) :=

theorem main_phase {c s : ℝ} (hc : 0 < c) (heps : ∀ j, S.eps j = epsCS c s j) :
    (s < 1 → IsEmpty (Stat S none)) ∧
      (s = 1 → c < 1 → Nonempty (Stat S none)) ∧
      (s = 1 → 1 ≤ c → IsEmpty (Stat S none)) ∧
      (1 < s → Nonempty (Stat S none)) := by
\end{lstlisting}

Item \emph{(2)} is three: the limit $C$, the tail asymptotic and the polynomial rate in one
statement; then the two senses in which the boundary data determine $C$ --- that they determine
the whole ladder, and that $C$ is a non-negative linear combination of them, with coefficients
the statement asserts to exist rather than exhibits.
\begin{lstlisting}[style=lean]
theorem main_invariant_measure {c : ℝ} (hc0 : 0 < c) (hc1 : c < 1)
    (heps : ∀ j, S.eps j = epsCS c 1 j) (L : Stat S none) :
    ∃ p : ℝ, psi c p = 0 ∧ 1 < p ∧ ∃ C : ℝ, 0 < C ∧
      Tendsto (fun m : ℕ => L.lam (.lad m) * (m : ℝ) ^ p) atTop (𝓝 C) ∧
      Tendsto (fun m : ℕ => (p - 1) * (m : ℝ) ^ (p - 1) * L.tailMass m) atTop (𝓝 C) ∧
      ∃ ϑ c₆ : ℝ, 0 < ϑ ∧ ϑ < 1 ∧ 0 < c₆ ∧
        ∀ m : ℕ, 1 ≤ m → |L.lam (.lad m) * (m : ℝ) ^ p - C| ≤ c₆ * (m : ℝ) ^ (-ϑ) := by
\end{lstlisting}
\begin{lstlisting}[style=lean]
theorem main_constant_determined {c : ℝ} (hc0 : 0 < c) (hc1 : c < 1)
    (heps : ∀ j, S.eps j = epsCS c 1 j) (L L' : Stat S none)
    (hagree : ∀ j : ℕ, 1 ≤ j → j ≤ S.d → L.lam (.lad j) = L'.lam (.lad j)) :
    ∀ m : ℕ, 1 ≤ m → L.lam (.lad m) = L'.lam (.lad m) := by

theorem main_constant_functional {c : ℝ} (hc0 : 0 < c) (hc1 : c < 1)
    (heps : ∀ j, S.eps j = epsCS c 1 j) :
    ∃ p : ℝ, psi c p = 0 ∧ 1 < p ∧ ∃ ν : ℕ → ℝ, (∀ j, 0 ≤ ν j) ∧
      0 < ∑ j ∈ Finset.Icc 1 S.d, ν j ∧ (∀ j, 1 ≤ j → 2 * j ≤ S.d → ν j = 0) ∧
      ∀ L : Stat S none,
        Tendsto (fun m : ℕ => L.lam (.lad m) * (m : ℝ) ^ p) atTop
          (𝓝 (∑ j ∈ Finset.Icc 1 S.d, ν j * L.lam (.lad j))) := by
\end{lstlisting}

Item \emph{(3)} is six. The first three are the unboundedness, at the finite exponents, at
$p=+\infty$, and at $p=2$ where it becomes $\widehat B=+\infty$, a divergent operator series
and a divergent mixing sum.
\begin{lstlisting}[style=lean]
theorem main_unbounded {c s : ℝ} (hc : 0 < c) (hs : 0 ≤ s)
    (heps : ∀ j, S.eps j = epsCS c s j) (L : Stat S none) {p : ℝ} (hp : 1 ≤ p) :
    (∀ t : ℝ, 0 < t → ∃ m, S.d < m ∧
        L.mass p (fun x => L.centredTail m x - pstar S none (L.centredTail m) x)
          ≤ t * L.mass p (L.centredTail m)) ∧
      ¬ ∃ B : ℝ, ∀ f : St → ℝ, (∃ C, ∀ x, |f x| ≤ C) → (∑' x, L.lam x * f x = 0) →
          L.mass p f ≤ B * L.mass p (fun x => f x - pstar S none f x) := by

theorem main_unbounded_infty {c s : ℝ} (hc : 0 < c)
    (heps : ∀ j, S.eps j = epsCS c s j) (L : Stat S none) :
    ¬ ∃ B : ℝ, ∀ (f : St → ℝ) (M Dd : ℝ),
        (∃ C, ∀ x, |f x| ≤ C) → (∑' x, L.lam x * f x = 0) → (∃ x, M ≤ |f x|) →
        (∀ x, |f x - pstar S none f x| ≤ Dd) → M ≤ B * Dd := by
\end{lstlisting}
\begin{lstlisting}[style=lean]
theorem main_unbounded_two {c s : ℝ} (hc : 0 < c) (hs : 0 ≤ s)
    (heps : ∀ j, S.eps j = epsCS c s j) (L : Stat S none) :
    (¬ ∃ Sop : Lp ℝ 2 L.mu →L[ℝ] Lp ℝ 2 L.mu,
        Sop * (1 - L.pstarL2 (rowOnChain_none S)) = 1 - L.piL2) ∧
      (¬ ∃ U : Lp ℝ 2 L.mu →L[ℝ] Lp ℝ 2 L.mu,
        Tendsto (partialSum (L.pstarL2 (rowOnChain_none S)) L.piL2) atTop (𝓝 U)) ∧
      ¬ Summable (L.betaHat (rowOnChain_none S)) := by
\end{lstlisting}

The other three are the backward length and the rate. $\texttt{main\_sigmaBar\_eq}$ gives the
closed form of item \emph{(3)}; $\texttt{main\_sigmaBar\_le}$ read at
$\texttt{cap}=\texttt{none}$ is what makes that length finite in the first place, and its other
half --- the same bound uniformly over every truncation --- is more than the theorem states; and
$\texttt{main\_rate}$ is the $\sqrt m$ rate \eqref{eq:lean_rate} along the trial functions.
\begin{lstlisting}[style=lean]
theorem main_sigmaBar_eq {c : ℝ} (hc0 : 0 < c) (hc1 : c < 1)
    (heps : ∀ j, S.eps j = epsCS c 1 j) :
    ⨆ n, sbar S none n = S.jbar / (1 - c) :=

theorem main_sigmaBar_le {c : ℝ} (hc0 : 0 < c) (hc1 : c < 1) {cap : Option ℕ}
    (heps : ∀ j, S.eps j = epsCS c 1 j) (n : ℕ) :
    sbar S cap n ≤ S.jbar / (1 - c) :=

theorem main_rate {c : ℝ} (hc0 : 0 < c) (hc1 : c < 1)
    (heps : ∀ j, S.eps j = epsCS c 1 j) (L : Stat S none) :
    ∃ p : ℝ, psi c p = 0 ∧ 1 < p ∧ ∃ c₁ c₂ : ℝ, 0 < c₁ ∧ c₁ ≤ c₂ ∧
      ∃ m₂ : ℕ, S.d < m₂ ∧ ∀ m : ℕ, m₂ ≤ m →
        c₁ * (m : ℝ) * L.mass 2 (fun x => L.centredTail m x - pstar S none (L.centredTail m) x)
          ≤ 4 * c₂ * (p - 1) * L.mass 2 (L.centredTail m) := by
\end{lstlisting}

Item \emph{(4)} is one statement, carrying both halves: the range is dense and proper, and a
pair of probability densities lies outside it.
\begin{lstlisting}[style=lean]
theorem main_unsolvable {c s : ℝ} (hc : 0 < c) (hs : 0 ≤ s)
    (heps : ∀ j, S.eps j = epsCS c s j) (L : Stat S none) :
    (L.defectRange ≤ L.kerPi ∧ L.defectRange.topologicalClosure = L.kerPi ∧
        L.defectRange ≠ L.kerPi) ∧
      ∃ fi fe : St → ℝ,
        (∀ x, 0 ≤ fi x) ∧ (∀ x, 0 ≤ fe x) ∧
        MemLp fi 2 L.mu ∧ MemLp fe 2 L.mu ∧
        ∑' x, L.lam x * fi x = 1 ∧ ∑' x, L.lam x * fe x = 1 ∧
        ¬ ∃ f : St → ℝ, MemLp f 2 L.mu ∧
            ∀ᵐ x ∂L.mu, f x - pstar S none f x = fi x - fe x :=
\end{lstlisting}

Item \emph{(5)} is three: the truncation is irreducible, it carries exactly one invariant
probability and a bounded diffusion operator, and that operator's norm grows at least like
$\sqrt K$.
\begin{lstlisting}[style=lean]
theorem main_truncation_irreducible {K : ℕ} (hK : Even K) (hdK : S.d ≤ K)
    (P : PreStat S (some K)) :
    (∀ x y : St, OnChain (some K) y → Reach S (some K) x y) ∧
      (∀ ⦃y : St⦄, OnChain (some K) y → 0 < P.lam y) :=

theorem main_truncation_bhat {K : ℕ} (hK : Even K) (hdK : S.d ≤ K) :
    Nonempty (Stat S (some K)) ∧
      (∀ L L' : Stat S (some K), L.lam = L'.lam) ∧
      ∀ L : Stat S (some K), ∃ Sop : Lp ℝ 2 L.mu →L[ℝ] Lp ℝ 2 L.mu,
        (1 - L.pstarL2 (rowOnChain_some hdK)) * Sop = 1 - L.piL2 ∧
        Sop * (1 - L.pstarL2 (rowOnChain_some hdK)) = 1 - L.piL2 ∧
        L.piL2 * Sop = 0 ∧ Sop * L.piL2 = 0 :=
\end{lstlisting}
\begin{lstlisting}[style=lean]
theorem main_truncation_sqrtK {c : ℝ} (hc0 : 0 < c) (hc1 : c < 1)
    (heps : ∀ j, S.eps j = epsCS c 1 j) :
    ∃ c₈ : ℝ, 0 < c₈ ∧ ∃ K₀ : ℕ, ∀ K : ℕ, K₀ ≤ K →
      ∀ (L : Stat S (some K)) (hrow : RowOnChain S (some K))
        (Sop : Lp ℝ 2 L.mu →L[ℝ] Lp ℝ 2 L.mu),
        Sop * (1 - L.pstarL2 hrow) = 1 - L.piL2 → c₈ * (K : ℝ) ≤ ‖Sop‖ ^ 2 := by
\end{lstlisting}

\begin{table*}[t]
\centering
\caption{Each item of Theorem~\ref{theo:lean_main}, the declarations of
$\texttt{GFNBounds/Doubling/Main.lean}$ that carry it, and the modules of
$\texttt{GFNBounds/Doubling/}$ in which the mathematics is proved.}
\label{tab:lean_map}
\small
\begin{tabular}{@{}c p{0.40\linewidth} p{0.36\linewidth}@{}}
\hline
item & declarations & modules \\
\hline
\emph{(1)} & \texttt{main\_irreducible}, \texttt{main\_phase} &
\texttt{Irreducible}, \texttt{PhaseEmpty}, \texttt{PhaseExists}, \texttt{StatExists} \\
\emph{(2)} & \texttt{main\_invariant\_measure}, \texttt{main\_constant\_determined},
\texttt{main\_constant\_functional} &
\texttt{DecayNotation}, \texttt{Product}, \texttt{Sharp}, \texttt{SharpFull},
\texttt{ConstantFunctional} \\
\emph{(3)} & \texttt{main\_unbounded}, \texttt{main\_unbounded\_infty},
\texttt{main\_unbounded\_two}, \texttt{main\_sigmaBar\_eq}, \texttt{main\_sigmaBar\_le},
\texttt{main\_rate} &
\texttt{Unbounded}, \texttt{Ramp}, \texttt{UnboundedL2}, \texttt{Family}, \texttt{Product},
\texttt{Supersolution}, \texttt{Length} \\
\emph{(4)} & \texttt{main\_unsolvable} & \texttt{Unsolvable} \\
\emph{(5)} & \texttt{main\_truncation\_irreducible}, \texttt{main\_truncation\_bhat},
\texttt{main\_truncation\_sqrtK} &
\texttt{Irreducible}, \texttt{TruncationStat}, \texttt{OperatorL2}, \texttt{Truncation},
\texttt{Product}, \texttt{RayleighBridge} \\
\hline
\end{tabular}
\end{table*}

The library is Lean 4 with Mathlib pinned at the tag \texttt{v4.31.0}, revision
\texttt{fabf563a7c95}. It compiles with $\texttt{warningAsError}:=\texttt{true}$, which turns
Lean's own ``declaration uses \texttt{sorry}'' warning into a compile error, so
$\texttt{sorry}$-freeness is enforced by the compiler rather than by discipline; a permissive
scaffold library exists for work in progress and is never imported by the strict one, so
graduating a result is a file move. An axiom sweep over the strict library reports
$4{,}006$ declarations in $99$ modules, each depending only on $\texttt{propext}$,
$\texttt{Classical.choice}$ and $\texttt{Quot.sound}$ --- Lean's three standard axioms, and
nothing else --- and a second audit parses $\texttt{\#print axioms}$ for the $245$ declarations
of an explicit audit file and fails on $\texttt{sorryAx}$ or on any axiom outside that set. A
kernel replay then re-checks every declaration of the strict library from the compiled
\texttt{.olean} files, with Mathlib imported and trusted, so that the certificate rests on the
kernel and not on the elaborator that produced the environment. A traceability check maps each
statement of the three appendices the library covers --- the $41$ of
Appendix~\ref{app:doubling}, the $37$ of the appendix of proofs and the $6$ of
Appendix~\ref{sec:silva_comparison} --- to its Lean declarations and to a \texttt{sha256} of
that statement's own \LaTeX{} block, so that editing a statement marks it stale until the map
is re-affirmed; it reports $84$ statements mapped and every digest current. Coverage of the
long-form appendix this one replaces is $29$ statements closed, $9$ partial and $3$ open.
\emph{Partial} means that some items of a multi-part statement are proved and others are not:
Lemma~\ref{lem:doubling_operator} is carried in its doubling-graph instance alone,
Lemma~\ref{lem:doubling_expansion} and Remark~\ref{rem:doubling_parity} are blocked on an
Euler--Maclaurin estimate Mathlib does not have, and Lemma~\ref{lem:doubling_weight} is carried
in a consequence form; none of the nine touches a clause of Theorem~\ref{theo:lean_main}. The
open three are remarks that appendix does not prove either, and
\texttt{docs/COVERAGE.md} gives the per-statement account. The whole gate is one command,
\texttt{make check}, last run green on 2026-09-13.

What this buys is one half of the reading and not both. The compiler certifies that each proof
proves its statement; what is left to the reader is whether the statement is the paper's, and
that is a question about \S\ref{sec:lean_model} and about \S\ref{sec:lean_scope} --- which is
why the model is quoted in full above, and why the scope below is part of the claim rather than
an apology attached to it.

\subsection{The ideas}\label{sec:lean_ideas}

This subsection says what the proofs do. It proves nothing, and nothing below is used anywhere
else in the paper.

\paragraph{One identity carries the invariant measure.}
Stationarity across the cut $A=[m,+\infty)$, at any $m>d$, has exactly two terms. The only way
out of $A$ is the decrement from $m$, which carries the probability $1-\varepsilon(m)$; the
only way in is a doubling from the window $W(m):=\{\lceil m/2\rceil,\dots,m-1\}$, since $s_0$
goes only to $s_f$ and $s_f$ lands in $\{1,\dots,d\}$ with $d<m$. Equating the two fluxes,
\begin{equation}\label{eq:lean_cut}
    \lambda_m\bigl(1-\varepsilon(m)\bigr)\;=\;\sum_{j\in W(m)}\lambda_j\,\varepsilon(j) ,
\end{equation}
and everything the certificate proves about the invariant measure is this one identity. The
certificate takes \eqref{eq:lean_cut}, and not the chain, as its hypothesis: the decay results
are stated for an arbitrary positive sequence satisfying the balance, with no invariance, no
normalization and no recurrence. That is why they apply verbatim to $\lambda^K$ below the cap,
where the same two-term count holds for every $d<m\leq K/2$, and it is why the library needs no
Markov-chain theory at all to prove them.

\paragraph{Why the exponent solves $c(2^p-1)=p$.}
Doubling multiplies the index, so the natural variable on the ladder is $\log_2$ of the state
and \eqref{eq:lean_cut} is a Laplace transform of one dyadic step. Feed the power ansatz
$\lambda_j\approx A j^{-p}$ into the balance at $s=1$, where $\varepsilon(j)=c/(j+1)$: the left
side is $Am^{-p}(1+o(1))$, the window sum is a Riemann sum for
$cA\int_{m/2}^mu^{-p-1}\,du=cAm^{-p}(2^p-1)/p$, and dividing out $Am^{-p}$ leaves the
Cram\'er equation \eqref{eq:lean_cramer} and nothing else. Its non-zero root is $>1$ exactly
when $c<1$ and $>0$ exactly when $c<1/\ln2$, so the two thresholds of the phase diagram are the
two crossings of a single root.

\paragraph{From ansatz to theorem: a chain that descends the ladder.}
An ansatz is not a theorem, and what turns it into one is a change of variable. Rescale to
$u_j:=\lambda_j j^{p_*}$; then \eqref{eq:lean_cut} says that $u$ at a level is a weighted average
of its values on the window below, the weights summing not to $1$ but to $R_0(m)=1+O(1/m)$.
Read the normalized weights as the transition law of a Markov chain descending the ladder. One
of its steps falls by at least one and by at most half --- from $y$ it lands in
$\{\lceil y/2\rceil,\dots,y-1\}$ --- so the descent from $m$ reaches a fixed base block
$[\ell,2\ell)$ without overshooting it, in at most $m$ steps. That count is linear and carries
nothing; what makes the descent fast is that the halving holds \emph{in probability}. From $y$
the chain lands at or below $y/\sqrt2$ with probability at least a fixed $\gamma$, so the time
it spends in any one dyadic block has expectation at most $2/\gamma^2$, and there are only
$\log_2(m/\ell)$ blocks between the start and the base block. The profile at an index is the
expectation of the accumulated weight times the value of $u$ where the descent lands, and it is
that bounded expected sojourn per block --- not a bounded number of steps --- which lets the
$1+O(1/m)$ deviations telescope over a whole descent into a single bounded factor
$e^{\pm O(1/\ell)}$, the constant depending on $c$ alone. The profile at every index above the
base block is then squeezed between the extremes of the profile on that one finite block, which
is the two-sided decay $c_1 j^{-p_*}\leq\lambda_j\leq c_2 j^{-p_*}$ --- with the same pair
$0<c_1\leq c_2$ that \eqref{eq:lean_rate} binds --- and no asymptotic hypothesis anywhere.

\paragraph{From two constants to one: a coupling.}
The sandwich leaves the constant anywhere between $c_1$ and $c_2$, and closing that gap is a
mixing statement about the descent. Two descents started in one dyadic block leave it with laws
that overlap: the exit law charges at least a fixed fraction $\omega$ of one reference law on
the block beneath --- a Doeblin minorization --- so the two descents can be coupled, and the
total variation between their landing laws falls geometrically in the number of blocks crossed.
The sandwich becomes a limit: $\lambda_j j^{p_*}\rightarrow C$ for some $C>0$, with a
polynomial rate $\left|\lambda_j j^{p_*}-C\right|\leq c_6 j^{-\vartheta}$ for some $c_6>0$ and
the exponent $\vartheta=\ln\bigl(1/(1-\omega)\bigr)/\ln\bigl(2/(1-\omega)\bigr)$ that the
minorization produces.

Two flags belong in the text beside it. The exponent is explicit and \textbf{not effective}: at
$c=1/2$ the proved $\omega$ is $0.0016$ against a block-exit minorization measured at $0.18$ on
exact exit laws, so the proved $\vartheta$ is $0.0024$ against $0.22$, while the observed rate
is $O(1/m)$. And $C$ is a non-negative linear functional of the boundary data and no more:
\eqref{eq:lean_cut} is homogeneous, so no equation in $c$ alone can determine it.

\paragraph{One trial function makes the operator unbounded.}
Apply $\mathrm{Id}-P_\star$ to the indicator of a tail and almost everything cancels: what
survives is supported on $\{m\}\cup W(m)$ --- the mass leaving at $m$, the mass entering over
the window --- so $\|(\mathrm{Id}-P_\star)f_m\|^2_{L^2(\lambda)}\leq2\lambda_m$, while
$\|f_m\|^2_{L^2(\lambda)}=L(m)\bigl(1-L(m)\bigr)$ because $f_m$ is a centred indicator. The
Rayleigh ratio at the cut $m$ is therefore at most
$\sqrt{2\lambda_m/\bigl(L(m)(1-L(m))\bigr)}$, and the operator is unbounded as soon as
$L(m)/\lambda_m$ is unbounded --- as soon as the tail fails to be comparable to its own last
term.

It cannot be comparable. Were $L(m)/\lambda_m$ bounded by $Q$, the tail would fall geometrically
at the rate $e^{-1/Q}$; but iterating the doubling edges along the powers of two gives
$L(2^D)\geq\lambda_1\prod_{i<D}\varepsilon(2^i)$ at every integer $D\geq1$, and the growth
condition --- that $\sum_{i<D}\log\bigl(1/\varepsilon(2^i)\bigr)$ is $o(2^D)$ as
$D\rightarrow\infty$ --- makes that far too large. The mechanism is
geometric and not spectral: the doubling edges plant mass exponentially far out along the
ladder, and no policy whose doubling probability decays subexponentially can keep the tail
comparable to its last term. The condition indicts the family and not the graph --- a
geometrically decaying $\varepsilon(j)=a\rho^j$ fails it, and no boundedness is claimed for any
policy here. At $p=+\infty$ the sharp indicator gives nothing and clipped ramps take its place:
their height grows with the clipping level while their defect stays bounded by a constant.

\paragraph{From unbounded to unsolvable.}
Restricted to $\ker\Pi$, the operator $\mathrm{Id}-P_\star$ is injective with dense range ---
only the constants are fixed by $P_\star$, only the constants are fixed by $P$, and the two
facts are adjoint to each other --- but its Rayleigh infimum is zero, so it cannot also be
surjective: the open mapping theorem would hand back the bounded inverse that the infimum
forbids. Any $\psi$ of zero mean outside the range then splits into two probability densities
$\finit,\fterm$ whose difference is a multiple of $\psi$, and that pair cannot be balanced in
$L^2(\lambda)$ at all. Two flags again. The pair comes from non-surjectivity and is
\textbf{not explicit}; and because the range is \emph{dense}, weak universality --- a defect
infimum of zero --- is untouched, and remains open on the infinite chain.

\paragraph{The truncation restores a finite constant, and prices it.}
Deleting the doubling edges above an even cap leaves a finite irreducible chain, so $\lambda^K$
exists, is unique and is positive, and $\widehat B_K$ is finite. Finiteness becomes growth by
reading the same per-cut inequality at a cut near $K/4$ and feeding it the two-sided decay,
which \eqref{eq:lean_cut} transfers to $\lambda^K$ unchanged, and $\widehat B_K$ grows at
least like $\sqrt K$.
This is loose, and by exactly one factor of $\sqrt K$ --- what \S\ref{sec:doubling_cost}
measures is $K$.

\paragraph{And the backward length stays finite.}
At $s=1$ the drift of the backward chain is $c-1$ at every ladder state, independent of the
state --- $s=1$ is the only exponent at which it is --- so the linear function $j\mapsto j/(1-c)$
solves the hitting-time recursion exactly, which bounds $\mathbb E(\sigma\mid X_0=j)$ above. The
matching lower bound is the one place where formalizing changed the proof.
Appendix~\ref{app:doubling} reaches it by optional stopping on a martingale, dominated by the
integrable excursion maximum; the certificate builds no martingale and never consumes the
excursion bound. It iterates the ladder operator on the identity $x\mapsto x$ instead --- that
operator sends the identity to $x\mapsto x-(1-c)$, so the martingale identity becomes an
induction on the iterate --- and it replaces dominated convergence by one explicit
supersolution, $W(m)=\sum_{i\leq m}\bigl(\lfloor\log_2i\rfloor+1\bigr)$, which the ladder
operator raises by at most $2$ and which grows like $m\log_2m$, so that the identity is
dominated by an arbitrarily small multiple of it plus a constant. The iterates of the identity therefore vanish, which is
$\mathbb E(\sigma\mid X_0=j)=j/(1-c)$; averaging over the target row gives
$\bar\sigma=\bar\jmath/(1-c)$. This is the divorce the appendix exists for: the condition of
\citet{morozov2025revisiting} is repaired at every $c<1$, and $\widehat B=+\infty$ at every one
of them.

\subsection{What the certificate does not carry}\label{sec:lean_scope}

Six clauses and conventions of Theorem~\ref{theo:doubling_main} are absent from
Theorem~\ref{theo:lean_main}, and they are the price of the rest being certified.

The first is uniqueness of the invariant probability on the loop closure. Positive recurrence
in Theorem~\ref{theo:doubling_main}\emph{(1)} carries it; $\texttt{main\_phase}$ delivers
existence alone, and the library discloses that it does not prove uniqueness there. At a finite
cap uniqueness is proved --- it is half of item \emph{(5)} --- by an argument that does not
transfer to the infinite chain.

The words \emph{transient} and \emph{null recurrent} of the phase diagram, and the clause
$\mathbb E(\sigma\mid X_0=j)=+\infty$ at $1\leq c<2$, are not stated. Mathlib \texttt{v4.31.0}
carries no discrete-time Markov chain theory and the library builds none, so each row of item
\emph{(1)} reads instead as the existence or the non-existence of an invariant probability,
which is what the rest of the appendix uses it for; Proposition~\ref{prop:doubling_phase}
carries the labels. The count of $K+2$ states of the truncation is not stated either: the Lean
state type is infinite and the truncation is carried by a predicate on it. The functional
$\bar\sigma$ is the supremum of the truncated hitting expectations, defined by the paper's own
recursion, rather than an expectation over a chain. And Appendix~\ref{app:doubling}'s three
unproved remarks --- the renewal reading of the balance, the absence of a log-periodic term,
and the second-order expansion --- are not formalized, that appendix not proving them either.

Four dependence clauses of constants go the same way. Theorem~\ref{theo:doubling_main}\emph{(3)}
says that its $c_7$ depends only on $c$, $d$ and an initial block of the invariant measure, and
Theorem~\ref{theo:doubling_sharp} says as much of its $c_6$; item \emph{(2)} above binds the
$\nu_j$ existentially, where they are in fact determined by $c$ and $d$ alone; and item
\emph{(5)} binds $c_8$ and $K_0$ the same way, where
Theorem~\ref{theo:doubling_main}\emph{(5)} names what they depend on. None of the four is
carried by the quoted existentials. They are carried instead, as explicit
formulas in $c$, $d$ and the initial block of the profile, by the library's explicit-constant
lemmas $\texttt{Decay.decay\_two\_sided\_explicit}$ and $\texttt{Decay.sharp\_explicit}$, and
for the $\nu_j$ by $\texttt{Decay.constant\_functional}$, whose coefficients depend on $c$ and
$d$ alone.

Two constants also change shape between the two statements, which is worth naming because a
reader comparing them will trip on it: Theorem~\ref{theo:doubling_main}\emph{(3)} writes its
rate as $\|f_m\|\geq c_7\sqrt m\,\|(\mathrm{Id}-P_\star)f_m\|$ where \eqref{eq:lean_rate}
carries the pair $0<c_1\leq c_2$ from which that $c_7$ is assembled, and it writes
$\widehat B_K\geq c_8\sqrt K$ where item \emph{(5)} states the same bound squared, so that the
$c_8$ of item \emph{(5)} is the square of the $c_8$ of
Theorem~\ref{theo:doubling_main}\emph{(5)}.

One clause of the certificate drops a hypothesis: the $\sqrt K$ lower bound is stated for every
$K\geq K_0$, even or not. The drop is empty. At an odd $K\geq K_0$ the truncation carries no
invariant probability at all, so the statement quantifies over nothing --- the state $K$ is on
chain and would have to carry positive mass, while its only sources are a decrement from $K+1$,
which is off chain, a doubling from $K/2$, which is not an integer, and the target row, which is
supported in $\{1,\dots,d\}$ while $K\geq K_0>d$. So this is a vacuous generalization and not a
strengthening, and item \emph{(5)} states the even case, which is the one the setting defines.

Where Appendix~\ref{app:doubling} says a thing is open, formal or non-explicit, the certificate
says so too, rather than laundering it.

The library is included as anonymized supplementary material together with its audit scripts;
\texttt{make check} builds it, runs the axiom sweep and the axiom audit, and re-derives the
per-statement map that ties each of the $84$ statements of the three appendices the library
covers to the declarations certifying it.

